% registry: KAKUNO
\methodhead{KAKUNO}{Kakuno and Kawada (1979): octahedral criterion with amplitude and mean terms}\label{sheet:KAKUNO}

\block{Source} H.~Kakuno, Y.~Kawada, \emph{Fatigue Engng Mater.\ Struct.} 2 (1979) 229--236
\cite{kakuno1979}, Eqs.~(4), (6)--(12); Figs.~2--6.

\block{Idea and scope} A correction of Sines' criterion: the octahedral shear stress amplitude at
the fatigue limit decreases linearly with the octahedral normal stress, with \emph{different}
slopes for its amplitude and its mean. Sines keeps only the mean term, which forces
$\tm/\fm=1/\sqrt3$ without mean stresses. Proposed for fatigue limits of round bars under repeated
bending and torsion with superimposed static bending and torsion; plain and moderately notched
members.

\block{Inputs}
\begin{center}\small
\begin{tabular}{@{}llp{100mm}@{}}
\toprule
Symbol & Unit & Meaning\\\midrule
$\sigma_{a_i}$, $\sigma_{m_i}$ ($i=1,2,3$) & MPa & principal stress amplitudes and principal mean stresses\\
$\sigma_a,\sigma_m$; $\tau_a,\tau_m$ & MPa & bending and torsion amplitude and mean (round bar)\\
$\sigma_{wb}=\fm$, $\tau_w=\tm$ & MPa & fully reversed bending and torsion limits\\
$\sigma_{up}=f_0$ & MPa & pulsating ($R=0$) bending limit, maximum stress of the cycle\\
$k_1=\sigma_{up}/(2\sigma_{wb})$ & -- & ratio of the pulsating amplitude to the reversed limit\\
$A$, $B$, $C$ & MPa, --, -- & material constants, Eqs.~\eqref{eq:kk-A}--\eqref{eq:kk-C}\\
\bottomrule
\end{tabular}
\end{center}

\block{Formulation} The criterion (Eq.~4) is
\begin{equation}
\tfrac13\sqrt{(\sigma_{a_1}-\sigma_{a_2})^2+(\sigma_{a_2}-\sigma_{a_3})^2+(\sigma_{a_3}-\sigma_{a_1})^2}
 = A - B\,(\sigma_{a_1}+\sigma_{a_2}+\sigma_{a_3}) - C\,(\sigma_{m_1}+\sigma_{m_2}+\sigma_{m_3}) ,
\label{eq:kk}
\end{equation}
whose left-hand side is the octahedral shear stress amplitude $\tau_{oct,a}=\sqrt{2/3}\sqrt{J_{2,a}}$.
For bending and torsion it becomes $\tfrac13(2\sigma_a^2+6\tau_a^2)^{1/2}=A-B\sigma_a-C\sigma_m$
(Eq.~6). Imposing the reversed torsion limit, the reversed bending limit and the pulsating bending
limit ($\sigma_a=\sigma_m=\sigma_{up}/2$) gives (Eqs.~7--9)
\begin{align}
A &= \sqrt{\tfrac23}\,\tau_w , \label{eq:kk-A}\\
B &= \frac{\sqrt2}{3}\left(\sqrt3\,\frac{\tau_w}{\sigma_{wb}} - 1\right), \label{eq:kk-B}\\
C &= \sqrt{\tfrac23}\,\frac{1-k_1}{k_1}\,\frac{\tau_w}{\sigma_{wb}} . \label{eq:kk-C}
\end{align}
Substituting, the limit curve of combined bending and torsion (Eq.~10) is
\begin{equation}
\Bigl(\frac{\tau_a}{\tau_w}\Bigr)^2\frac{1}{p^2}
+ \Bigl(\frac{\sigma_a}{\sigma_{wb}}\Bigr)^2\frac{\tfrac{2}{\sqrt3}\tfrac{\sigma_{wb}}{\tau_w}-1}{p^2}
+ \frac{\sigma_a}{\sigma_{wb}}\,\frac{2-\tfrac{2}{\sqrt3}\tfrac{\sigma_{wb}}{\tau_w}}{p} = 1 ,
\qquad p = 1 - \frac{\sigma_m}{\sigma_{wb}}\frac{1-k_1}{k_1} ,
\end{equation}
and $p=1$ without mean stresses (Eq.~11). A torsional mean stress does not enter. The paper
recommends $k_1=0.83$ as an average for structural carbon and alloy steels.

\block{Constants and calibration} $A$, $B$, $C$ from $\tau_w$, $\sigma_{wb}$, $\sigma_{up}$; nothing
is fitted.

\block{Life} The paper gives none; finite lives use the master curve on the index below
(Section~\ref{sec:mastercurve}), a repository extension.

\block{Repository choices}
\begin{itemize}
\item The model returns the index
  $\bigl(\tau_{oct,a} + B\,\lvert\tr\boldsymbol\sigma_a\rvert + C\,\tr\Sm\bigr)/A$, equal to 1 on the
  limit, with $\boldsymbol\sigma_a$ the amplitude tensor; the sum of the principal amplitudes is
  taken as $\lvert\tr\boldsymbol\sigma_a\rvert$, positive like any amplitude.
\item Proportional loading only (principal stress amplitudes are undefined otherwise): a case with
  non-parallel $\Ss$, $\Sc$ is out of scope. The mean may have any direction.
\item $\sigma_{wb}$, $\tau_w$, $\sigma_{up}$ are the workbook $\fm$, $\tm$, $f_0$. Without $f_0$, $C$
  is undefined: a case with $\tr\Sm\neq0$ is out of scope; $k_1=0.83$ is not used as a default.
\end{itemize}

\block{Worked example} Ni--Cr--Mo steel of Fig.~5, stresses divided by $\sigma_{wb}$:
$\tau_w=0.617$, $k_1=0.83$, so $A=0.5038$, $B=0.0324$, $C=0.1032$. With $\sigma_m=0.46$,
$p = 1-0.46\cdot0.17/0.83 = 0.906$, and the curve meets the $\tau$ axis at
$\tau_a = p\,\tau_w = 0.559$: $\tau_{oct,a}=\sqrt{2/3}\cdot0.559=0.4563$, index
$(0.4563+0+0.1032\cdot0.46)/0.5038 = 1.000$, with or without a static torsion.

\block{Verification} Eqs.~(7)--(10) and (12) hold to machine precision. Fig.~5: the $\tau$
intercept is 0.559 against 0.557 read from the figure; the $\sigma$ intercept is 0.906 against
about 0.93 drawn (unexplained). The paper prints no table.
